% !PRINTED 96
\weberpage{96}
%<<CONT
signs there enters in each case a complete system of residues of $p$ (with the exclusion of 0).

But we wish to remain with the assumption $i_1=3$, and to understand by $\varepsilon_1, \varepsilon_2, \ldots \varepsilon_\mu$ any combination of the numbers $\pm 1$. Then we can also write the congruence (3) thus
\weq{4}{\varepsilon_1\alpha_1+\varepsilon_2\alpha_2+\cdots+\varepsilon_\mu\alpha_\mu\equiv 0\ (\mathrm{mod}\ 3).}

The combination of signs $\varepsilon_1, \varepsilon_2, \ldots \varepsilon_\mu$, in which one sign can be assumed arbitrarily, since the simultaneous change of all the signs in (4) alters nothing, determines the individual group $\mathfrak{A}$, and evidently these groups $\mathfrak{A}$ are also all different from one another. For if, e.g., in one group $\varepsilon_1=\varepsilon_2=+1$, and in another $\varepsilon_1=+1$, $\varepsilon_2=-1$, then the first contains the combination of numbers:
\[
\alpha_1\equiv 1,\quad \alpha_2\equiv -1,\quad \alpha_3\equiv 0,\quad \ldots\ \alpha_\mu\equiv 0\ (\mathrm{mod}\ 3),
\]
which does not occur in the second.

In order to find the groups $\mathfrak{A}$ definitively, one takes as a basis a system of primitive roots $g_1, g_2, \ldots g_\mu$ of $q_1, q_2, \ldots q_\mu$ and determines $a$ to the modulus $m$ from the congruences:
\weq{5}{\begin{aligned}
a &\equiv g_1^{\alpha_1}\ (\mathrm{mod}\ q_1)\\
&\equiv g_2^{\alpha_2}\ (\mathrm{mod}\ q_2)\\
&\cdots\cdots\cdots\cdots\cdots\\
&\equiv g_\mu^{\alpha_m}\ (\mathrm{mod}\ q_\mu). % [sic: Weber prints \alpha_m; the pattern requires \alpha_\mu.]
\end{aligned}}
If we set, understanding by $\sigma$ one of the residues 0, $\pm 1$,
\weq{6}{\sigma \equiv \varepsilon_1\nu_1+\varepsilon_2\nu_2+\cdots+\varepsilon_\mu\nu_\mu\ (\mathrm{mod}\ 3),}
then $\mathfrak{A}$ consists of all numbers $a$ whose indices $\nu_1=\alpha_1, \nu_2=\alpha_2, \ldots \nu_\mu=\alpha_\mu$ satisfy the condition $\sigma=0$. The two other values $\sigma=1$, $\sigma=-1$ determine the two subsidiary groups
$\mathfrak{A}'$, $\mathfrak{A}''$ of $\mathfrak{A}$, so that
\[
\mathfrak{N} = \mathfrak{A} + \mathfrak{A}' + \mathfrak{A}''
\]
holds. If $a'$ and $a''$ are each a number from $\mathfrak{A}'$ and $\mathfrak{A}''$, then $\mathfrak{A}'=\mathfrak{A}a'$, $\mathfrak{A}''=\mathfrak{A}a''$.

We now denote by $\eta$ the cyclotomic period that belongs to $\mathfrak{A}$, and by $\eta'$, $\eta''$ the conjugate periods, thus
\weq{7}{\eta = \sum r^a,\quad \eta' = \sum r^{a'},\quad \eta'' = \sum r^{a''}.}

If then $\varrho$ is an imaginary third root of unity, we obtain the resolvent
\[
(\varrho, \eta) = \eta + \varrho\,\eta' + \varrho^2\,\eta'',
\]
%CONT>>
